% GENERATED FILE. Edit canonical lessons, then run npm run generate:books.
% canonical-source-hash: fnv1a64:97461c520d9f0d19
% canonical-lessons: PA-C191-amount-return, PA-C191-ghatt, PA-C191-ghatt-review, PA-C191-paani-ghatt, PA-C191-samaa-ghatt, PA-C191-thaan-ghatt, PA-C191-payoff

\chapter{Less Water, Time, and Space}
\label{ch:pa-a2-ghatt-quantity}
\hlchaptermodality{191}

\begin{chapteropening}
\textbf{By the end of this chapter:} \emph{I can use one familiar quantity word to say there is not much water, time, or space.}

Retrieve three limited-amount statements through listening, speaking, reading, and writing.
\end{chapteropening}

\section[\texorpdfstring{\pa{ਪਾਣੀ} · \pa{ਸਮਾਂ} · \pa{ਥਾਂ} (paani; samaa; thaan)}{ਪਾਣੀ · ਸਮਾਂ · ਥਾਂ (paani; samaa; thaan)}]{Return to water, time, and place}

\label{lesson:PA-C191-amount-return}



\begin{quote}
\emph{Hear:} \emph{paani; samaa; thaan}; YOU SAY: \textbf{\pa{ਪਾਣੀ}; \pa{ਸਮਾਂ}; \pa{ਥਾਂ}}
\end{quote}

\begin{grammarlens}[title={Grammar Lens}]
Keep three familiar nouns ready. One names water, one names time, and one names place. The next lesson adds one quantity word that can describe each noun.
\end{grammarlens}

\subsection*{Your turn}
\begin{itemize}
\raggedright
  \item \emph{Hear:} \emph{paani}; YOU SAY: \textbf{\pa{ਪਾਣੀ}}
  \item \emph{Hear:} \emph{samaa}; YOU SAY: \textbf{\pa{ਸਮਾਂ}}
  \item \emph{Hear:} \emph{thaan}; YOU SAY: \textbf{\pa{ਥਾਂ}}
\end{itemize}

\begin{tcolorbox}[breakable,colback=teal!4,colframe=teal!35!black,title={Before you move on}]
Say “water,” “time,” then “place.” This opener owns no new form.
\end{tcolorbox}
\section[\texorpdfstring{\pa{ਘੱਟ} (ghatt)}{ਘੱਟ (ghatt)}]{One word for a limited amount}

\label{lesson:PA-C191-ghatt}



\begin{quote}
Retrieve “water,” “time,” and “place” before describing their amount.
\end{quote}

\begin{grammarlens}[title={Grammar Lens: direct owner}]
\begin{quote}
\textbf{\pa{ਘੱਟ}} — \emph{ghaṭṭ} — little, less, or not much
\end{quote}

Use the bounded meaning “not much” here. No comparison word is needed. The next lesson retrieves this one form before it enters a complete sentence.
\end{grammarlens}

\begin{tcolorbox}[breakable,skin=enhanced,colback=violet!6,colframe=violet!45!black,boxrule=0.5pt,arc=1mm,left=6pt,right=6pt,top=4pt,bottom=4pt,fonttitle=\bfseries,title={Across the family}]
Turner groups Punjabi \emph{ghaṭṭ}, “diminished,” with Prakrit \emph{ghaṭṭaï}, “is ruined,” under reconstructed \emph{ghaṭṭati}, “decreases, is wanting.” The old idea and the modern quantity meaning stay close. Source: \href{https://dsal.uchicago.edu/cgi-bin/app/soas\_query.py?qs=\%2Agha\%E1\%B9\%AD\%E1\%B9\%ADati\&searchhws=yes}{Turner, CDIAL 4415}.
\end{tcolorbox}

\subsection*{Your turn}
\begin{itemize}
\raggedright
  \item \emph{Hear:} \emph{ghatt}; YOU SAY: \textbf{\pa{ਘੱਟ}}
  \item \emph{Read it:} \textbf{\pa{ਘੱਟ}}
  \item \emph{Write it:} copy \textbf{\pa{ਘੱਟ}} once
\end{itemize}

\begin{tcolorbox}[breakable,colback=teal!4,colframe=teal!35!black,title={Before you move on}]
Produce the word for “not much” without looking.
\end{tcolorbox}
\section[\texorpdfstring{\pa{ਘੱਟ} (ghatt)}{ਘੱਟ (ghatt)}]{Retrieve the quantity word}

\label{lesson:PA-C191-ghatt-review}



\begin{quote}
Hear “not much” and prepare the single Punjabi word.
\end{quote}

\begin{grammarlens}[title={Grammar Lens}]
Keep \textbf{\pa{ਘੱਟ}} bounded to “little, less, or not much.” This lesson adds no opposite word and no new comparison frame.
\end{grammarlens}

\subsection*{Your turn}
\begin{itemize}
\raggedright
  \item \emph{Hear:} \emph{ghatt}; YOU SAY: \textbf{\pa{ਘੱਟ}}
  \item \emph{Read it:} \textbf{\pa{ਘੱਟ}}
  \item \emph{Write it:} cover the model and write \textbf{\pa{ਘੱਟ}}
\end{itemize}

\begin{tcolorbox}[breakable,colback=teal!4,colframe=teal!35!black,title={Before you move on}]
Say the meaning, then produce the written form.
\end{tcolorbox}
\section[\texorpdfstring{\pa{ਪਾਣੀ} \pa{ਘੱਟ} \pa{ਹੈ} (paani ghatt hai)}{ਪਾਣੀ ਘੱਟ ਹੈ (paani ghatt hai)}]{Build the water frame}

\label{lesson:PA-C191-paani-ghatt}



\begin{quote}
Retrieve “water,” “not much,” and “is.”
\end{quote}

\begin{grammarlens}[title={Grammar Lens: first complete statement}]
\begin{quote}
\textbf{\pa{ਪਾਣੀ} \pa{ਘੱਟ} \pa{ਹੈ}} — \emph{pāṇī ghaṭṭ hai} — there is not much water
\end{quote}

Read the familiar noun, the new quantity word, and the familiar final verb as one small practical statement.
\end{grammarlens}

\subsection*{Your turn}
\begin{itemize}
\raggedright
  \item \emph{Hear:} \emph{paani ghatt hai}; YOU SAY: \textbf{\pa{ਪਾਣੀ} \pa{ਘੱਟ} \pa{ਹੈ}}
  \item \emph{Read it:} \textbf{\pa{ਪਾਣੀ} \pa{ਘੱਟ} \pa{ਹੈ}}
  \item \emph{Write it:} cover and write \textbf{\pa{ਪਾਣੀ} \pa{ਘੱਟ} \pa{ਹੈ}}
\end{itemize}

\begin{tcolorbox}[breakable,colback=teal!4,colframe=teal!35!black,title={Before you move on}]
Produce “there is not much water.”
\end{tcolorbox}
\section[\texorpdfstring{\pa{ਸਮਾਂ} \pa{ਘੱਟ} \pa{ਹੈ} (samaa ghatt hai)}{ਸਮਾਂ ਘੱਟ ਹੈ (samaa ghatt hai)}]{Build the time frame}

\label{lesson:PA-C191-samaa-ghatt}



\begin{quote}
Retrieve “time,” “not much,” and “is.”
\end{quote}

\begin{grammarlens}[title={Grammar Lens: complete statement}]
\begin{quote}
\textbf{\pa{ਸਮਾਂ} \pa{ਘੱਟ} \pa{ਹੈ}} — \emph{samā̃ ghaṭṭ hai} — there is not much time
\end{quote}

Keep the quantity frame from the water statement. Replace only the familiar noun at the front.
\end{grammarlens}

\subsection*{Your turn}
\begin{itemize}
\raggedright
  \item \emph{Hear:} \emph{samaa ghatt hai}; YOU SAY: \textbf{\pa{ਸਮਾਂ} \pa{ਘੱਟ} \pa{ਹੈ}}
  \item \emph{Read it:} \textbf{\pa{ਸਮਾਂ} \pa{ਘੱਟ} \pa{ਹੈ}}
  \item \emph{Write it:} cover and write \textbf{\pa{ਸਮਾਂ} \pa{ਘੱਟ} \pa{ਹੈ}}
\end{itemize}

\begin{tcolorbox}[breakable,colback=teal!4,colframe=teal!35!black,title={Before you move on}]
Produce “there is not much time.”
\end{tcolorbox}
\section[\texorpdfstring{\pa{ਥਾਂ} \pa{ਘੱਟ} \pa{ਹੈ} (thaan ghatt hai)}{ਥਾਂ ਘੱਟ ਹੈ (thaan ghatt hai)}]{Move the quantity word to space}

\label{lesson:PA-C191-thaan-ghatt}



\begin{quote}
Retrieve “place,” “not much,” and “is.”
\end{quote}

\begin{grammarlens}[title={Grammar Lens: second familiar noun}]
\begin{quote}
\textbf{\pa{ਥਾਂ} \pa{ਘੱਟ} \pa{ਹੈ}} — \emph{thā̃ ghaṭṭ hai} — there is not much space
\end{quote}

The quantity word stays unchanged. Only the familiar noun changes from time to place or space.
\end{grammarlens}

\subsection*{Your turn}
\begin{itemize}
\raggedright
  \item \emph{Hear:} \emph{thaan ghatt hai}; YOU SAY: \textbf{\pa{ਥਾਂ} \pa{ਘੱਟ} \pa{ਹੈ}}
  \item \emph{Read it:} \textbf{\pa{ਥਾਂ} \pa{ਘੱਟ} \pa{ਹੈ}}
  \item \emph{Write it:} cover and write \textbf{\pa{ਥਾਂ} \pa{ਘੱਟ} \pa{ਹੈ}}
\end{itemize}

\begin{tcolorbox}[breakable,colback=teal!4,colframe=teal!35!black,title={Before you move on}]
Produce “there is not much space.”
\end{tcolorbox}
\section[\texorpdfstring{\pa{ਪਾਣੀ} \pa{ਘੱਟ} \pa{ਹੈ} · \pa{ਸਮਾਂ} … (paani ghatt hai …)}{ਪਾਣੀ ਘੱਟ ਹੈ · ਸਮਾਂ … (paani ghatt hai …)}]{A bounded quantity payoff}

\label{lesson:PA-C191-payoff}



\begin{quote}
Recall the quantity word, then prepare the three familiar nouns.
\end{quote}

\subsection*{Your turn: four skills}
\begin{itemize}
\raggedright
  \item \emph{Hear:} \emph{paani ghatt hai; samaa ghatt hai; thaan ghatt hai}; YOU SAY: \textbf{\pa{ਪਾਣੀ} \pa{ਘੱਟ} \pa{ਹੈ}; \pa{ਸਮਾਂ} \pa{ਘੱਟ} \pa{ਹੈ}; \pa{ਥਾਂ} \pa{ਘੱਟ} \pa{ਹੈ}}
  \item \emph{Read it:} \textbf{\pa{ਪਾਣੀ} \pa{ਘੱਟ} \pa{ਹੈ} · \pa{ਸਮਾਂ} \pa{ਘੱਟ} \pa{ਹੈ} · \pa{ਥਾਂ} \pa{ਘੱਟ} \pa{ਹੈ}}
  \item \emph{Write it:} cover the line and write \textbf{\pa{ਘੱਟ}} once, then one complete statement
\end{itemize}

\begin{tcolorbox}[breakable,colback=teal!4,colframe=teal!35!black,title={Before you move on}]
Produce all three complete statements once more.

This is the fourteenth prerequisite tranche exposed by Punjabi A2 mock drafting. It does not unblock Punjabi A2 mock 1 by itself. It is not an A2 mock, not A2 readiness or attainment, and not external or human validation.
\end{tcolorbox}
